\documentclass[11pt]{article}
\usepackage{amsmath, amssymb}
\usepackage[margin=1in]{geometry}

\newcommand{\pit}{\pi_\theta}
\newcommand{\piref}{\pi_{\text{ref}}}
\newcommand{\given}{\,\vert\,}
\newcommand{\E}{\mathbb{E}}
\newcommand{\KL}{\mathrm{KL}}

\title{Direct Preference Optimization --- Derivation}
\author{Eshan Kemp}
\date{SF Phase, rep 1}

\begin{document}
\maketitle

%======================================================
\section{The RLHF objective and its optimal policy}
%======================================================

The constrained RLHF objective: maximize reward, penalize divergence
from the reference policy.
\begin{equation}
  \max_{\pit}\ \E_{x,\, y \sim \pit(\cdot \given x)}
  \big[ r(x,y) \big]
  \;-\; \beta\, \KL\!\big( \pit(y \given x) \,\|\, \piref(y \given x) \big)
  \label{eq:objective}
\end{equation}

\paragraph{Goal.} Show the optimal policy is Boltzmann over the reference:
\begin{equation}
  \pi^*(y \given x) \;=\; \frac{1}{Z(x)}\,
  \piref(y \given x)\, \exp\!\Big( \tfrac{1}{\beta}\, r(x,y) \Big),
  \qquad
  Z(x) = \sum_{y} \piref(y \given x)\, \exp\!\Big( \tfrac{1}{\beta}\, r(x,y) \Big).
  \label{eq:optimal}
\end{equation}

\paragraph{Derivation.} % expand the KL, rewrite the objective as a single
% expectation, recognize it as a negative KL to an (unnormalized) target,
% conclude the optimum.

% --- your work here ---

\bigskip
\noindent\textbf{Physics bridge.} % r = negative energy, beta = temperature,
% piref = density of states / prior, Z(x) = partition function.
% Note the beta -> 0 (greedy / zero-temperature) and beta -> infty
% (reverts to reference) limits and what they mean for the policy.

% --- your note here ---

%======================================================
\section{Solving for the implicit reward}
%======================================================

Invert \eqref{eq:optimal} to express the reward in terms of the policies
and the (intractable) partition function:
\begin{equation}
  r(x,y) \;=\;
  \beta \log \frac{\pi^*(y \given x)}{\piref(y \given x)}
  \;+\; \beta \log Z(x).
  \label{eq:reward}
\end{equation}

% --- show the inversion ---

%======================================================
\section{The partition function cancels (the trick)}
%======================================================

Bradley--Terry preference model over a chosen/rejected pair
$(y_c, y_r)$ at a \emph{shared} prompt $x$:
\begin{equation}
  P(y_c \succ y_r \given x)
  = \sigma\!\big( r(x, y_c) - r(x, y_r) \big),
  \qquad \sigma(z) = \frac{1}{1+e^{-z}}.
  \label{eq:bt}
\end{equation}

\paragraph{Claim.} Substituting \eqref{eq:reward} into the reward
\emph{difference} $r(x,y_c) - r(x,y_r)$, the $\beta \log Z(x)$ terms cancel.

% --- show the cancellation. Same x => same Z(x) => gone.
% Physics: Z never survives a ratio of Boltzmann factors at one temperature. ---

\begin{equation}
  r(x,y_c) - r(x,y_r)
  = \beta \log \frac{\pi^*(y_c \given x)}{\piref(y_c \given x)}
  - \beta \log \frac{\pi^*(y_r \given x)}{\piref(y_r \given x)}.
  \label{eq:rewarddiff}
\end{equation}

%======================================================
\section{The DPO loss}
%======================================================

Replace $\pi^*$ with the trainable $\pit$, take the negative log-likelihood
of the Bradley--Terry model over the preference dataset:
\begin{equation}
  \mathcal{L}_{\text{DPO}}(\theta)
  = -\,\E_{(x, y_c, y_r)}
  \left[
    \log \sigma\!\left(
      \beta \log \frac{\pit(y_c \given x)}{\piref(y_c \given x)}
      - \beta \log \frac{\pit(y_r \given x)}{\piref(y_r \given x)}
    \right)
  \right].
  \label{eq:dpoloss}
\end{equation}

% --- connect each of the four log-terms to a get_logprobs() call in code ---

%======================================================
\section{Implementation bridge: response log-probability}
%======================================================

A response log-probability is a sum of per-token log-probs under
teacher forcing --- over \emph{response} tokens only, never the prompt:
\begin{equation}
  \log \pi(y \given x)
  = \sum_{t=1}^{|y|} \log \pi\big( y_t \given x, y_{<t} \big).
  \label{eq:logprob}
\end{equation}

% --- mark explicitly: t indexes response tokens; prompt tokens are
% conditioning context, excluded from the sum. THIS is the equation the
% prompt-masking in get_logprobs() enforces. ---

\end{document}